\documentclass[11pt]{article}
\usepackage[a4paper,margin=26mm]{geometry}
\usepackage{amsmath,amssymb,amsthm,booktabs,hyperref}
\newtheorem{theorem}{Theorem}\newtheorem{lemma}[theorem]{Lemma}
\newcommand{\F}{\mathbb F}\newcommand{\gen}[1]{\langle#1\rangle}
\title{Three explicit unitary examples in rank four}
\author{A certificate for the examples in Section 7 of Leemans--Vandenschrick}\date{}
\begin{document}\maketitle
The source reports three pairs of enantiomorphic chiral polytopes for
$\operatorname{PSU}_3(17)$, of types $[3,7,3]$, $[3,8,3]$ and $[3,9,3]$.
It gives neither matrices nor a search certificate for this sentence.
We supply explicit matrices and verify these three positive examples.
We do not infer that other types are absent, and we do not treat the
source's historical question for larger fields as currently open.

Let $F=\F_{17}[\beta]$, where $\beta^2=3$. Since $3$ is a nonsquare
modulo $17$, the nontrivial field automorphism is
$\overline\beta=-\beta$. We use the standard Hermitian form
$z^*z=\overline z^{\,t}z$. For a nonisotropic column $z$, put
\[
 R_z=-I+\frac{2zz^*}{z^*z}.
\]
This matrix fixes $z$, negates its orthogonal complement, and has
unitary determinant one. Put
\[
 u=R_{(4\beta,0,1)^t},\qquad w=R_{(0,0,1)^t}.
\]
For the three choices below, set $v=R_{(x,y,1)^t}$ and
\[
                 a=wv,\qquad b=vwu,\qquad c=uw.       \tag{1}
\]
\[
\begin{array}{c|c|c|c}
 |b|&x&y&\operatorname{Tr}_{F/\F_{17}}(x/(4\beta))\\\hline
 7&3+4\beta&1+3\beta&2\\
 8&7+\beta&1+3\beta&9\\
 9&8+6\beta&4+\beta&3
\end{array}                                                   \tag{2}
\]
Here and below orders are projective orders when the matrices represent
unitary projective transformations.

\begin{theorem}
For each row of (2), the images of $(a,b,c)$ generate
$\operatorname{PSU}_3(17)$ and satisfy the full string rotation
intersection property. They define a rank-four chiral polytope of the
indicated type $[3,|b|,3]$. The three types give three distinct
pairs of enantiomorphic markings.
\end{theorem}
The proof combines exact, small-input finite group calculations with a
full-automorphism argument. We retain the standard rotation
correspondence for abstract polytopes, the classical order and faithful
isotropic-point action of $\operatorname{PSU}_3(17)$, and its classical
automorphism description
\[
 \operatorname{Aut}(\operatorname{PSU}_3(17))
       =\operatorname{PGU}_3(17)\rtimes\gen{\overline{\phantom{x}}}.
                                                               \tag{3}
\]
For this prime field the last factor has order two. The inverse-transpose
automorphism on unitary matrices is the same field map, so no separate
graph coset is omitted. These are attributed classical group inputs,
not new classifications asserted here. GAP's exact finite permutation
group algorithms supply the group orders and subgroup intersections
specified below; neither the source search nor a named-group constructor
supplies them.

\section*{The interval relations and actual subgroup checks}
All four nontrivial centres in (2) and the definition of $u$ have
Hermitian norm four. In particular all denominators are nonzero.
The first and last rotations have order three. One may check this
directly, or use the formula for two unitary involutions: their product
has eigenvalue one and trace $4\lvert z^*e_3\rvert^2/(z^*z)-1$.
Here that trace is zero, so the remaining two eigenvalues are primitive
cube roots. Equation (1) gives identically
\[
                       ab=u,\qquad bc=v,\qquad abc=w. \tag{4}
\]
The other required interval products are therefore involutions.
Exact multiplication gives the middle orders in (2).

The finite calculation constructs all $4,914$ isotropic projective
points directly. They are the normalized vectors
\[
 (1,s,t)\quad\hbox{with }1+s\bar s+t\bar t=0,
 \qquad (0,1,s)\quad\hbox{with }1+s\bar s=0.
\]
The remaining normalized point $(0,0,1)$ is not isotropic. These
lists contain every projective point of norm zero exactly once.
The program uses row representatives and right multiplication, preserving
these Hermitian equations; the abstract matrix products are unchanged.
The action of each literal matrix in (1) on these lists is computed,
without identifying it with a preselected permutation generator.
All three generated permutation groups have order
\[
 \frac{17^3(17^3+1)(17^2-1)}{3}=2,317,678,272.
\]
They are subgroups of the faithful projective unitary action; equality
with its classical order proves generation of the whole group.
The two adjacent subgroups have respective orders
\[
\begin{array}{c|r|r}
 |b|&|\gen{a,b}|&|\gen{b,c}|\\\hline
 7&168&168\\
 8&2448&2448\\
 9&2448&2448
\end{array}
\]
No group identification from these last orders is needed.

We check every indexed intersection, rather than infer it from these
orders. For $I\subseteq\{0,1,2,3\}$ let $G_I$ be generated by those
of the following six elements whose endpoint pair lies in $I$:
\[
 (01):a,\ (02):ab,\ (03):abc,\ (12):b,\ (13):bc,\ (23):c.
\]
The verifier constructs all sixteen actual permutation subgroups and
checks $G_I\cap G_J=G_{I\cap J}$ for all $256$ ordered pairs.
For nested index sets it verifies actual subgroup containment; for
all other pairs it computes the actual intersection. Adding the
trivial endpoint conventions in the source does not change any of
these groups. The three complete tests pass, giving precisely the
rotation intersection property.

\section*{Excluding every reversing automorphism}
\begin{lemma}
Let $u,w$ be the two matrices above. If $v=R_{(x,y,1)^t}$ and
$\operatorname{Tr}_{F/\F_{17}}(x/(4\beta))\ne0$, no automorphism
of $\operatorname{PSU}_3(17)$ fixes $u,w$ and sends $v$ to $wvw$.
\end{lemma}
\begin{proof}
Write $\alpha=4\beta$. A linear matrix commuting with $w$ preserves
its centre $e_3$ and its axis $\gen{e_1,e_2}$. Commuting also with
$u$, whose centre has nonzero first and third coordinates, forces
its two off-diagonal entries in the $e_1,e_2$ block to vanish and
its first and third diagonal entries to agree. Thus its projective
form is $D_r=\operatorname{diag}(1,r,1)$. It is a unitary similitude
exactly when $r\bar r=1$.

Fixing the projective involutions really does give these commuting
matrix conditions: their determinant-one involutory lifts are unique,
since the scalar centre has odd order three. Conjugation by a projective
unitary representative, and the field map, preserve these lifts.
The semilinear map given by field conjugation followed by
$D_0=\operatorname{diag}(\alpha/\bar\alpha,1,1)$ fixes $u,w$.
Composing with it reduces every semilinear centralizer element to the
linear case. Hence the full centralizer in (3) consists of the eighteen
$D_r$ and the eighteen maps $D_rD_0\overline{\phantom{x}}$.

The centre of $wvw$ is $(-x,-y,1)^t$. Its last coordinate is already
normalized, so equality of centres under a linear centralizer element
requires $x=-x$, hence $x=0$. Under a semilinear element it requires
$(\alpha/\bar\alpha)\bar x=-x$, or
\[
              x/\alpha+\overline{x/\alpha}=0.
\]
Both possibilities contradict the hypothesis. This accounts for every
automorphism in (3), independently of the second coordinate equation.
\end{proof}

A reversing automorphism for (1) would send $a$ to $a^{-1}$ while
fixing $ab=u$ and $abc=w$, by the standard chiral criterion.
Since $v=wa$, it would send $v$ to $wa^{-1}=wvw$.
The three traces in (2) are nonzero. The lemma therefore excludes all
reversing automorphisms. With generation, (4) and the verified indexed
intersections, the rotation correspondence proves the theorem.
Distinct displayed types distinguish the three enantiomorphic pairs.

\section*{Certificate and scope}
The supplied GAP program builds the field, the three homologies and all
isotropic points afresh. It checks group membership, all exact rotation
orders, all $768$ indexed intersections and the $36$ possible reversers
for each row. An independent elementary field implementation checks the
matrix identities, types and traces without importing GAP.
Two additional controls prevent weaker tests from being mistaken for
the theorem. The centre $(\beta,1+2\beta,1)$ gives the same type
$[3,8,3]$ but only a group of order $2448$, so fails generation.
The centre $(1,6,1)$ has trace zero and admits an actual reversing
semilinear map, so cannot pass the chirality test.

The exploratory normalization and sixteen completed intermediate cases
are retained with their receipts. They are not an asserted census of
all ranks or all types. In particular, our statement recovers the
source's three reported positive pairs and does not turn the phrase
``a computer search gave also three pairs'' into an unproved assertion
that there are no others. The positive claim is attributed to
Leemans--Vandenschrick; the explicit centres and centralizer certificate
are supplied here. No new question resolution, Lean proof, or numerical
compression measurement is claimed.
\end{document}
